\title{Parameter-Efficient Orthogonal Finetuning via Factorization}

\begin{abstract}
Large-scale foundation models have become increasingly prevalent, yet their immense computational and financial costs render full retraining infeasible. Consequently, methods for cost-effective model adaptation to specific tasks are crucial. In this work, we investigate a systematic approach—Orthogonal Finetuning (OFT)—for customizing foundation models to downstream applications. Although OFT offers strong generalization capabilities, its parameter requirements remain substantial, primarily due to the inherent dimensionality of orthogonal matrices involved. To enhance parameter efficiency, we analyze OFT through the lens of information transfer and identify essential criteria for more efficient representations. Drawing inspiration from the Cooley-Tukey fast Fourier transform, which efficiently transmits information, we introduce a compact orthogonal parameterization based on butterfly architectures. This leads us to develop Orthogonal Butterfly (BOFT), a novel, parameter-efficient adaptation scheme grounded in orthogonal structure. Notably, BOFT generalizes OFT, treating it as a particular instance within a broader framework. Comprehensive experiments adapting large-scale vision transformers, language models, and text-to-image diffusion architectures to a variety of vision and language tasks demonstrate the effectiveness of our approach.
\end{abstract}

\section{Introduction}

Breakthroughs in models such as ChatGPT~\cite{brown2020language,chen2021evaluating} and Stable Diffusion~\cite{rombach2022high} have highlighted the extraordinary generalization achieved by modern foundation models. However, these advances have come with a surge in model size, often involving billions of parameters

(\emph{e.g.}, GPT-3 encompasses approximately 175B parameters), presenting substantial obstacles to training such models from the ground up. Consequently, progress in the area now hinges on adapting pre-existing models, obviating the need for expensive retraining. That is, efficient strategies for customizing foundation models to downstream tasks are imperative. There are three main adaptation paradigms: \emph{model finetuning}~\cite{hulora2022,zaken2022bitfit,guo2020parameter,raffel2020exploring,qiu2023controlling,zhang2022adaptive,chavan2023one}, which leaves the model architecture unaltered while modifying a subset of weights; \emph{adapter tuning}~\cite{rebuffi2017learning,houlsby2019parameter,pfeiffer2020adapterfusion,lin2020exploring,he2021towards}, which introduces small trainable modules into the network; and \emph{prompt tuning}~\cite{li2021prefix,lester2021power}, where trainable prompts are prepended to model inputs. Among these, model finetuning stands out for its straightforward implementation and effectiveness, in addition to maintaining unchanged inference latency.

Model finetuning is grounded in the idea that the updated (finetuned) model should remain closely aligned with its pretrained counterpart according to specific metrics, thereby retaining the knowledge acquired during pretraining. Typically, prevalent finetuning strategies employ a reduced learning rate as a heuristic, implicitly minimizing the divergence between the pretrained and finetuned parameters. Recognizing that weight matrices possess inherent structure, some methodologies advocate for the preservation of the relational properties within these matrices—specifically, maintaining consistent pairwise angles among neurons. Building on this observation, the Orthogonal Finetuning (OFT) framework~\cite{qiu2023controlling} was introduced, which enforces that all neurons within a layer undergo the same orthogonal transformation, rigorously conserving pairwise neuron angles throughout training. OFT has shown notable improvements in both generalization and convergence for applications such as text-to-image diffusion model finetuning~\cite{qiu2023controlling}. Nevertheless, the orthogonal matrices in OFT can become very large, leading to a substantial number of parameters to optimize. To alleviate this, OFT proposes the use of block-diagonal orthogonal matrices, which restricts parameter count. Yet, this reduction comes with significant trade-offs—the fixed block-diagonal pattern imposes strict structural sparsity and treats blocks independently, which may limit the representation capacity (for instance, precluding the approximation of more general linear mappings) and could induce suboptimal inductive biases. Critically, the question of how to optimally select a block sparsity pattern for orthogonal transformations is still unresolved.

A central challenge in this context is how to construct a dense orthogonal matrix in a manner that minimizes parameter usage. Although at first it might appear unmanageable, given that a dense orthogonal matrix of size $d$ requires $\mathcal{O}(d^2)$ parameters, our approach departs from conventional methods by expressing a dense orthogonal matrix as a product of multiple structured sparse matrices. This method relies on the principle that parameter count can be curtailed by substituting additional computation for memory usage. Since our orthogonal matrix arises from the successive multiplication of sparse matrices, these products must be computed in every training step. To conceptualize this factorization, we draw an analogy to information transmission: the synthesis of a dense orthogonal matrix through sparse matrix products can be mapped to the transmission of data across a grid-like graph structure. This analogy supports a flexible framework, allowing diverse designs for sparse matrix factorizations that both restrict the parameter budget and retain sufficient expressivity to yield dense matrices.

Within this information transmission paradigm, we achieve parameter efficiency by leveraging butterfly network structures, as found in the Cooley-Tukey fast Fourier transform algorithm~\cite{cooley1965algorithm}, renowned for their efficient communication between nodes~\cite{klasing1994broadcasting}. Notably, the butterfly network in Cooley-Tukey suggests a method of sparse matrix factorization referred to as butterfly factorization. Employing this strategy, we can synthesize a $d$-dimensional dense matrix by multiplying $\mathcal{O}(\log d)$ sparse matrices, each possessing $\mathcal{O}(d)$ non-zero elements. When every constituent sparse matrix is constrained to be orthogonal, this yields a parameterization with merely $\mathcal{O}(d\log d)$ parameters—a drastic reduction compared to the standard approach. Building upon this efficient construction, we introduce Orthogonal Butterfly (BOFT), a novel finetuning strategy that achieves parameter efficiency by parameterizing with sparse orthogonal matrices in the butterfly pattern. BOFT generalizes OFT as a specific instance, providing a versatile and orthogonally structured finetuning framework. The block-diagonal and butterfly constructions share the crucial property of introducing sparsity, which decreases the number of parameters that must be learned. In comparison, LoRA~\cite{hulora2022} relies on low-rank approximations; both low-rank and sparse matrices constitute important classes of structured matrices~\cite{candes2011robust} enabling parameter-efficient model adaptation.

In contrast to the block-diagonal formulation utilized in OFT, which balances between expressiveness and regularization, BOFT leverages the butterfly pattern to enable a more gradual transition between matrices belonging to the full orthogonal group (maximal expressiveness) and identity matrices (maximal regularity). This design facilitates the identification of a more compact hypothesis space within the orthogonal group that is suitable for downstream applications. The butterfly structure is widely employed in efficient linear transformations, including the discrete Fourier and discrete cosine transforms. Building on this, we posit that our structured parameter-efficient approach inherently introduces inductive biases, which promote generalizability and reduce susceptibility to overfitting. This perspective is grounded in the observation that while the butterfly architecture can replicate a broad array of classical linear transforms, such capacity is unattainable for the block-diagonal structures adopted by OFT. Our main contributions are as follows:

\begin{itemize}[leftmargin=*,nosep]
    \item We address the challenge of achieving parameter efficiency in orthogonal finetuning by introducing a novel framework based on information transmission. This formulation recasts the design of parameter-efficient, dense orthogonal matrices as an information flow problem over grid-like graphs.
    \item Drawing inspiration from the butterfly networks utilized in the Cooley-Tukey algorithm, we develop the Orthogonal Butterfly method—a parameter-efficient technique for orthogonal finetuning—formulated within the context of the information transmission framework.
    \item We present several theoretical arguments elucidating why BOFT succeeds in drastically reducing the number of tunable parameters while preserving high expressiveness and stability during training. The factorization inherent in BOFT also naturally supports weight interpolation (refer to Figure~\ref{fig:boft_interpolation}).
    \item To our knowledge, this is the first application of orthogonal finetuning beyond controllable text-to-image generation~\cite{qiu2023controlling}, demonstrating the broad utility of the approach as a versatile method for model finetuning. We evaluate BOFT on a diverse set of adaptation problems, spanning domains from computer vision to natural language processing, where it consistently surpasses leading contemporary methods, thereby substantiating its parameter-efficiency and generalization capability.
\end{itemize}

\section{Related Work}

\textbf{Parameter-efficient finetuning (PEFT)}. With the ever-growing scale and capability of foundation models (\emph{e.g.}, \cite{brown2020language,radford2021learning,rombach2022high,kirillov2023segment}), optimizing these models for specific downstream applications without altering most of their parameters has attracted significant research attention~\cite{treviso2023efficient,ding2023parameter,lialin2023scaling}. Within the diverse landscape of PEFT strategies~\cite{houlsby2019parameter,aghajanyan2020intrinsic,hulora2022,edalati2022krona,wang2022adamix,gheini2021cross,zaken2022bitfit,guo2020parameter,sung2021training,ansell2022composable,lester2021power,li2021prefix,vu2022spot,he2021towards,mao2021unipelt,karimi2021compacter,liu2022few,sung2022lst,chen2022parameter,jia2022visual,chen2022adaptformer,zhang2022neural,jie2023fact,lian2022scaling,luo2023towards}, methods based on reparameterization~\cite{aghajanyan2020intrinsic,hulora2022,edalati2022krona} are particularly aligned with the direction of our study. The LoRA approach~\cite{hulora2022} modifies the pretrained weights by incorporating the product of two low-rank matrices, achieving effective results in a variety of natural language tasks. While LoRA fixes the rank across all added matrices, recent approaches \cite{zhang2022adaptive,valipour2022dylora,zhang2023increlora} permit the rank to vary across layers to optimize parameter distribution. Due to its conceptual clarity, such low-rank reparameterizations are widely adopted~\cite{dettmers2023qlora,zi2023delta,chavan2023one}. Informed by findings on the role of hyperspherical energy in model generalization~\cite{liu2018learning,liu2021orthogonal}, \cite{qiu2023controlling} introduced orthogonal finetuning as an alternative mechanism to adapt text-to-image diffusion models. Orthogonal finetuning (OFT) entails learning an orthogonal transformation for the neurons within a given layer, delivering not only improved generalization but also more stable optimization dynamics relative to LoRA. Notwithstanding these benefits, OFT typically involves a greater number of learnable parameters compared to LoRA, prompting the need for more efficient versions. Furthermore, the suitability of OFT beyond the setting of text-to-image diffusion model adaptation~\cite{qiu2023controlling} remains untested. BOFT builds upon OFT by utilizing butterfly factorization to reduce parameter count, making it more efficient. This advancement makes it feasible to leverage the strengths of orthogonal finetuning for a broader set of adaptation scenarios.

\textbf{Butterfly structures}. The radix-2 Cooley-Tukey algorithm decomposes the $N$-point discrete Fourier transform into a sequence of $\frac{N}{2}$-point transforms, thereby giving rise to the characteristic butterfly structure, which can be formalized as a sequence (product) of sparse matrices—collectively known as a butterfly matrix. Such butterfly matrices have found application as a means for parametrizing orthogonal transformations to sidestep pivoting in Gaussian elimination, thereby facilitating more efficient computation~\cite{parker1995random}. Additional usage scenarios include bolstering the training stability of recurrent neural networks~\cite{jing2017tunable}, and enabling more efficient kernel approximations~\cite{munkhoeva2018quadrature}. Efforts by \cite{dao2019learning,dao2019kaleidoscope} further show how fast linear transformations may be achieved through butterfly-based parametrization, while \cite{chen2022pixelated,dao2022monarch} illustrate their benefit for the efficient training of neural models. Moreover, butterfly structures prove valuable in rapid matrix-vector operations~\cite{michielssen1996multilevel,o2010algorithm}, approximating data-sparse matrices~\cite{li2015butterfly}, and enhancing network transmission~\cite{klasing1994broadcasting,li2003linear,rajasingh2016transmission}. Unlike existing approaches, our emphasis is on leveraging butterfly structures specifically to boost the parameter efficiency of OFT.

\section{An Information Transmission View on Orthogonal Finetuning}

We begin by reviewing fundamental concepts of OFT. In OFT, rather than directly modifying the pretrained weight matrix, the approach reparameterizes the adapted weights as the product of a trainable orthogonal matrix and the fixed pretrained matrix. Unlike LoRA, which adjusts weights through the addition of a low-rank matrix, OFT applies a multiplicative update using an orthogonal transformation. To maintain parameter-efficiency, LoRA employs a low-rank decomposition, whereas standard OFT adopts a block-diagonal orthogonal matrix as its structure~\cite{qiu2023controlling}; the efficiency increases as the size of each diagonal block decreases. Figure~\ref{fig:comp_lora} provides a visual comparison between these strategies. The central idea behind utilizing an orthogonal transformation for finetuning is to retain pairwise angular relationships between neurons~\cite{liu2018learning,liu2021orthogonal,qiu2023controlling}, thereby largely preserving the semantic information acquired during pretraining. Specifically, OFT seeks to optimize an orthogonal matrix $\bm{R}\in\mathbb{R}^{d\times d}$ corresponding to a pretrained linear layer with weight matrix $\bm{W}^0\in\mathbb{R}^{d\times n}$, modifying the standard forward computation from $\bm{z}=(\bm{W}^0)^\top\bm{x}$ to $\bm{z}=(\bm{R}\bm{W}^0)^\top\bm{x}$, with $\bm{x}\in\mathbb{R}^{d}$ as input and $\bm{z}\in\mathbb{R}^{n}$ as output. To ensure the adaptation process begins from the original pretrained model, the matrix $\bm{R}$ is initialized as the identity. To preserve its orthogonality during finetuning, we employ the Cayley parameterization method, following \cite{qiu2023controlling,liu2021orthogonal}: $\bm{R}=(\bm{I}+\bm{Q})(\bm{I}-\bm{Q})^{-1}$, where $\bm{Q}$ denotes a skew-symmetric matrix satisfying $\bm{Q}=-\bm{Q}^\top$. For parameter-efficiency, a block-diagonal design is applied by representing the orthogonal matrix as $\text{diag}(\bm{R}_1,\bm{R}_2,\cdots,\bm{R}_r)$, where each block $\bm{R}_i\in\mathbb{R}^{b\times b}$ for $i=1, \ldots, r$, and the total dimension $br=d$. This efficiency, however, comes with a trade-off: the block-diagonal setup presupposes that each neuron’s dimensions (i.e., columns of $\bm{W}^0$) are partitioned into $r$ groups, and that each group undergoes independent orthogonal transformation. Although the block-diagonal design demonstrates empirical benefits, the rationale for splitting neuron dimensions solely based on their indices is questionable, suggesting a preference for dense orthogonal matrices. This leads to a natural inquiry: \emph{Is it possible to design a dense orthogonal matrix while maintaining parameter-efficiency?}

To investigate this issue, we suggest constructing a dense orthogonal matrix by multiplying several sparse orthogonal matrices together. To offer a consistent and clear framework for analyzing the sparsity structure within orthogonal matrix factorization, we conceptualize the process of forming a dense orthogonal matrix in OFT as analogous to an information transmission system. More precisely, creating a dense matrix $\bm{R}\in\mathbb{R}^{d\times d}$ as the product of $m$ square matrices, $\thickmuskip=2mu \medmuskip=2mu \bm{R} = \bm{B}_m\bm{B}_{m-1}\cdots\bm{B}_1$, can be mapped to an information transfer across a grid of $d\times (m+1)$ nodes, as depicted in Figure~\ref{fig:info_trans}. The rationale for this analogy originates from recognizing that a dense $d$-by-$d$ matrix embodies a full connectivity pattern from one group of $d$ nodes to another. In the matrix $\bm{R}$, an entry $R_{ij} = 0$ signifies that no information passes from node $j$ to node $i$, while a nonzero $R_{ij}$ implies that such information flow is possible. Therefore, the expression of $\bm{R}$ as a sequence of matrices $\bm{B}_m\bm{B}_{m-1}\cdots\bm{B}_1$ can equivalently be seen as orchestrating a multi-stage information transfer, where each matrix $\bm{B}_i$ induces its own underlying graph. The transfer of information sequentially follows the pathway set by $\bm{B}_1$ up through $\bm{B}_n$. As an illustrative case in Figure~\ref{fig:info_trans}, we analyze the decomposition $\bm{R} = \bm{B}_5\bm{B}_4\bm{B}_3\bm{B}_2\bm{B}_1$, showing both the sparsity patterns and the related graph representation. The graph in Figure~\ref{fig:info_trans} arises from unfolding the matrix multiplication across layers. Here, $\bm{B}_i$ functions as the adjacency matrix defining connectivity from the $i$-th layer’s nodes to those in layer $(i+1)$. Specifically, the entry $(j_1, j_2)$ in $\bm{B}_i$ encodes whether a directed link exists from node $j_2$ at the $i$-th layer to node $j_1$ at layer $(i+1)$ (with a zero entry indicating absence of a link). To ensure $\bm{B}_5\bm{B}_4\bm{B}_3\bm{B}_2\bm{B}_1$ results in a fully dense matrix, it is necessary for information originating at each first-layer node to reach every node in the $6$-th layer. Should we consider the product $\bm{R} = \bm{B}_4\bm{B}_3\bm{B}_2\bm{B}_1$, the path from first-layer nodes to fifth-layer nodes, there exists the case where, for example, information from node $1$ cannot reach node $3$, revealing a zero entry at position $(3,1)$ in the sparsity pattern of the product. This pattern recurs for smaller compositions, such as $\bm{R} = \bm{B}_3\bm{B}_2\bm{B}_1$ and $\bm{R} = \bm{B}_2\bm{B}_1$, consistently linking the induced graphs and the sparsity configurations. In summary, to produce a dense matrix $\bm{R}\in\mathbb{R}^{d\times d}$, the set of $m$ factor matrices $\bm{B}_m,\ldots,\bm{B}_1$ must together yield a collection of directed edges over a $d \times (m+1)$ node grid, where each edge only connects nodes at consecutive layers (i.e., columns), guaranteeing that information from any first-level node can be transferred to all nodes in the $(m+1)$-th layer.

Viewing OFT through the lens of information transmission offers valuable insight into the distribution of nonzero elements in its factorization matrices. As depicted in Figure~\ref{fig:blk_diag}, the original OFT yields a matrix $\bm{R}$ with a block-diagonal pattern. While this structure minimizes the number of trainable parameters, it is inherently incapable of forming a dense $\bm{R}$. The objective, therefore, is to construct a dense orthogonal matrix by combining $m$ sparse orthogonal factors, while utilizing the smallest possible number of trainable parameters. From the perspective of information transmission, two main principles emerge: \emph{(i)} \textbf{dense connectivity}, which requires that each node in the initial layer be linked, via some path, to every node in the final layer; and \emph{(ii)} \textbf{minimum free edges}, meaning the total number of edges should be minimized as allowed by the orthogonality condition. 

It is important to highlight that orthogonality imposes a stringent restriction on connections between consecutive layers. For instance, each factor $\bm{B}_i$ must be full-rank—an essential aspect of orthogonality—which necessitates $d$ edges that establish a bijection between the $i$-th and $(i+1)$-th levels. Consequently, there must be at least $d$ such connections (as illustrated, for instance, in Figure~\ref{fig:info_trans}, where $d = 4$). These specific $d$ edges are dictated by the requirement for orthogonality and do not count towards the trainable edge total since their associated values are fixed and not trainable (for example, in a $d\times d$ orthogonal matrix where only $d$ entries are nonzero, they must be strictly $\pm 1$). Due to the fewer number of parameters needed by orthogonal matrices compared to dense matrices, imposing the orthogonality condition creates dependencies among the presence of edges. As a case in point, within a $2\times 2$ orthogonal matrix, if the $(1,1)$ position is zero, orthogonality enforces that the $(2,2)$ position must also be zero—so omitting one edge can enforce the omission of another. Thus, for each admissible connection pattern, orthogonality may either necessitate the addition or removal of certain edges. In this context, examining the sparsity structure of the resulting orthogonal matrix through the information transmission lens is particularly illuminating for OFT, because the emphasis is on identifying which elements of $\bm{R}$ are trainable and nonzero, rather than on their exact values.

A straightforward approach to establishing dense connectivity between two layers requires $\mathcal{O}(d^2)$ connections, corresponding to a fully dense orthogonal matrix with $d^2-d$ edges (for instance, when $d=4$, this amounts to 12 connections). Figure~\ref{fig:info_trans} illustrates an example of a practical matrix factorization, which utilizes only $10$ connections overall—fewer than those used in a dense orthogonal matrix. This method allows for an examination of the parameter-efficiency of OFT through a graphical lens, facilitating the design of feasible factorizations with ease. The conceptual underpinning is inspired by the butterfly graphs employed in the Cooley-Tukey algorithm~\cite{cooley1965algorithm}, which are known for their efficient dense connectivity between $d$ sources and $d$ receivers using just $\mathcal{O}(d\log d)$ edges. For example, the configuration depicted in Figure~\ref{fig:info_trans} achieves dense connections with $10$ edges, yet a butterfly network requires only $8$. The discussion that follows details how the butterfly topology contributes to enhanced parameter efficiency.

\section{Orthogonal Parameterization by Butterfly Factorization}

Originating from the Cooley-Tukey algorithm for the fast Fourier transform, the butterfly structure exploits the property that a small alteration in the frequency domain leads to a widespread impact in the spatial domain. This property parallels the information transmission scenario considered here, where each node at the initial level can relay information to every node in the terminal level. Furthermore, butterfly topologies are widely adopted in computer networks for streamlined information dissemination~\cite{solihin2015fundamentals,li2011linear}. Let $k \geq 2$ be a power of $2$; we introduce the \emph{butterfly factor} $\bm{B}^F(k)$ as follows:

\begin{equation}\label{eq:but_factor}
\begin{aligned}
    \bm{B}^F(k)=\begin{bmatrix}
    \text{diag}(\bm{d}_1) & \text{diag}(\bm{d}_2) \\
    \text{diag}(\bm{d}_3) & \text{diag}(\bm{d}_4)
    \end{bmatrix}\in\mathbb{R}^{k\times k},
\end{aligned}
\end{equation}

where each vector $\bm{d}_i \in \mathbb{R}^{\frac{k}{2}}$ for $i=1,2,3,4$. For $d=2^N$, the $d$-dimensional \emph{butterfly matrix} $\bm{B}(d)\in\mathbb{R}^{d\times d}$ is defined recursively by

\begin{equation}
\begin{aligned}
    \!\!\bm{B}(d)=\tilde{\bm{B}}(d,d)\!\cdot\!\begin{bmatrix}
    \bm{B}_1(\frac{d}{2})\!\!\!\!\! & \bm{0} \\
    \bm{0}\!\!\!\!\! &\bm{B}_2(\frac{d}{2})
    \end{bmatrix}= \tilde{\bm{B}}(d,d)\tilde{\bm{B}}(d,\frac{d}{2})\cdots\tilde{\bm{B}}(d,2),
\end{aligned}
\end{equation}

where $\bm{B}_1(\frac{d}{2})$ and $\bm{B}_2(\frac{d}{2})$ represent two butterfly matrices, each of dimension $\frac{d}{2}$. Next, we introduce the \emph{butterfly component} defined by $\thickmuskip=2mu \medmuskip=2mu \tilde{\bm{B}}(d,k)=\text{diag}(\bm{B}^F_1(k),\cdots,\bm{B}^F_{d/k}(k))$, which constitutes a $d\times d$ block-diagonal matrix whose blocks each have size $k$; here, $\bm{B}^F_i(k)$ for all $i$ are the butterfly factors as stated in Equation~\ref{eq:but_factor}. This framework enables us to construct orthogonal matrices using butterfly matrices for parameterization. The key requirement is that every multiplicative factor $\tilde{\bm{B}}(d,k)$ within the butterfly matrix $\bm{B}(d)$ maintains orthogonality. Focusing first on the case where the block size is $2$, $\tilde{\bm{B}}(d,2)$, each block can be parametrized to be orthogonal by means of the Cayley transform (or equivalently as a $2$-dimensional rotation), consistent with approaches found in \cite{liu2021orthogonal,qiu2023controlling}. Since the arrangement of non-zero entries in each butterfly component is merely a permutation of that in $\tilde{\bm{B}}(d,2)$, all butterfly components can be readily structured as orthogonal matrices. As a result, this yields a resource-efficient representation of orthogonal matrices using a large number of $2\times 2$ orthogonal matrices as building blocks. Extending upon the methodology of \cite{chen2022pixelated}, we further generalize to a block butterfly component $\tilde{\bm{B}}^b(d,k)$, in which every entry in each vector $\bm{d}_i$ is substituted by a $b\times b$ matrix. To ensure the orthogonality of $\tilde{\bm{B}}^b(d,2)$, we parametrize each $2b\times 2b$ block as an orthogonal matrix. The support pattern of other butterfly components, specifically $\tilde{\bm{B}}^b(d,k)$ for $k>2$, are realized as block-wise permutations of the nonzero structure of $\tilde{\bm{B}}^b(d,2)$, making it straightforward to maintain orthogonality for all such blocks. In summary, assembling all these components yields the BOFT forward computation as

\begin{equation}
    \begin{aligned}\nonumber
    \bm{z}=\big(\bm{R}(m,b)\cdot\bm{W}^0\big)^\top\bm{x},~~~~\text{s.t.}~\left\{\bm{R}(m,b)=\prod_{i=1}^m \tilde{\bm{B}}^b_{(d,i)}~~\&~~\underbrace{\big(\tilde{\bm{B}}^b_{(d,j)}\big)^\top\tilde{\bm{B}}^b_{(d,j)}=\tilde{\bm{B}}^b_{(d,j)}\big(\tilde{\bm{B}}^b_{(d,j)}\big)^\top\!=\bm{I}_d}_{\forall j\in [1, m]}\right\},
    \end{aligned}
\end{equation}

Here, $\tilde{\bm{B}}^b_{(d,i)}$ stands in for $\tilde{\bm{B}}^b(d,2^{m - i+1})$ for ease of notation, and $\bm{I}_d$ indicates a $d \times d$ identity matrix. The orthogonal operator $\bm{R}(m,b)\in\mathbb{R}^{d\times d}$ is structured as the product of several orthogonal butterfly matrices. For simplicity, we let BOFT parametrized by $\bm{R}(m,\frac{b}{2})$ be referred to as BOFT($m$,$b$) where $b\geq 2$. When $\thickmuskip=2mu \medmuskip=2mu m=1$, BOFT($1$,$b$) simplifies to the block-diagonal OFT~\cite{qiu2023controlling} using a block size of $b$. Moreover, BOFT($1$,$d$) is equivalent to the original OFT~\cite{qiu2023controlling} equipped with an unrestricted orthogonal matrix. BOFT($\log \frac{2d}{b}$,$b$) constructs $\bm{R}$ as the block butterfly matrix $\bm{B}^{\frac{b}{2}}(d)$, resulting in a full dense orthogonal matrix. Generally, the parameterization of BOFT($m$,$b$) introduces $\frac{1}{2}(b-1)dm$ effective learnable parameters for adaptation of a linear transformation of dimensions $d\times n$. Setting $m=\log d,~b=2$ yields the butterfly configuration for BOFT, which scales with $\mathcal{O}(d\log d)$ parameters. In comparison, the standard OFT using a dense orthogonal matrix incurs $\mathcal{O}(d^2)$ parameter cost, while the block-diagonal OFT with $r$ blocks entails $\mathcal{O}(bd)$. Thus, whereas the original OFT necessitates setting $b=d$ for achieving a dense orthogonal transformation, BOFT affords flexibility in choosing $b$ to achieve density.

\textbf{Identity Initialization in BOFT}. In most finetuning protocols, the initialization aligns exactly with the weights of the original pretrained model, ensuring that modifications remain close to the original parameter regime. For instance, LoRA initializes its low-rank updates at zero. In the context of BOFT, we perform initialization by assigning identity matrices to all butterfly factors (meaning the skew-symmetric matrices employed in the Cayley transformation begin as zero matrices).

\textbf{BOFT Multiplicative Dropout}. To further alleviate overfitting, LoRA~\cite{hulora2022} integrates a Dropout layer into the low-rank modification, following the approach from standard Dropout~\cite{srivastava2014dropout}, which fits naturally for LoRA's additive structure. However, such dropout is not readily applicable to BOFT due to its inherently multiplicative updates. To bridge this gap, we introduce a multiplicative Dropout specifically designed for BOFT. Given that the orthogonal operator $\bm{R}(m,b)$ consists of $m$ butterfly factors, each of which can be rearranged as $2b \times 2b$ block-diagonal orthogonal blocks, multiplicative Dropout operates by (i) stochastically selecting $p_1\%$ of the butterfly matrices and, within each, (ii) choosing $p_2\%$ of their block-diagonal parts to be replaced with identity matrices.

\section{Intriguing Insights and Discussions}

\textbf{Expressivity of BOFT}. The combination of the butterfly structure with permutations has the capacity to exactly represent numerous well-known fast linear transformations~\cite{dao2019learning,dao2019kaleidoscope} (\emph{e.g.}, fast Fourier transform, Hadamard transform). However, the degree to which our orthogonal butterfly matrix is able to approximate an arbitrary orthogonal matrix has not been thoroughly explored. To address this, we perform a simulation experiment where we attempt to approximate a randomly generated dense orthogonal matrix~\cite{anderson1987generation} of dimensions $512\times 512$. In Figure~\ref{fig:para_eff}, we present the results averaged across 10 independent initializations. The vertical axis shows the approximation error, whereas the horizontal axis tracks the number of effective trainable parameters. Each curve of the same color indicates BOFT instances configured with the same block size; specifically, the point furthest to the left on a given curve corresponds to BOFT($1$,$b$) (\emph{i.e.}, the canonical block-diagonal OFT where each block has size $b$). The findings demonstrate that BOFT is generally more parameter-efficient compared to OFT. As a case in point, BOFT($9$,$2$) achieves better expressiveness than BOFT($1$,$16$), yet does so with a significantly smaller parameter count. More broadly, smaller block sizes $b$ combined with higher multiplicities $m$ confer greater parameter efficiency in BOFT; for example, BOFT($6$,$4$) employs far fewer parameters but achieves a comparable error to BOFT($2$,$16$). This suggests that the butterfly matrix parametrizes a more structured subset within the orthogonal group than the block-diagonal architecture, making BOFT theoretically more expressive than OFT with equivalent block sizes.

\begin{theorem}[Expressivity of BOFT]\label{thm:exp_boft}
    BOFT possesses greater expressive capacity than OFT when the block sizes are matched. Specifically, to represent every orthogonal matrix of dimension $d$ using butterfly matrices, one can form products of the form $\bm{B}_{d-1,1}(d)\bm{B}^\top_{d-1,2}(d)\cdots\bm{B}_{1,1}(d)\bm{B}^\top_{1,2}(d)$, with each $\bm{B}_{i,j}(d)$ representing a butterfly matrix for all admissible $i, j$.
\end{theorem}

Building on Theorem~\ref{thm:exp_boft}, we can consider a straightforward extension of BOFT where the final orthogonal matrix is constructed as $\thickmuskip=2mu \medmuskip=2mu \bm{R}^G(m_1,b_1,m_2,b_2,l)=\bm{R}_{l,1}(m_1,b_1)\bm{R}_{l,2}^T(m_2,b_2)\cdots\bm{R}_{1,1}(m_1,b_1)\bm{R}_{1,2}^T(m_2,b_2)$, in which $\bm{R}_{i,j}^T(m,b)$ denotes the orthogonal matrix instance used within BOFT. Notably, when $m_1=m_2=\log d$, $b_1=b_2=2$, and $l=d-1$, the generalized construction $\bm{R}^G(m_1,b_1,m_2,b_2,l)$ is capable of expressing the entirety of the orthogonal group—a hierarchy known in literature as the kaleidoscope hierarchy~\cite{dao2019kaleidoscope}. Nonetheless, it is important to recognize that increased expressivity does not guarantee superior finetuning performance: for example, full finetuning methods, while universally expressive, often produce suboptimal outcomes. Thus, finding an appropriate balance between model expressivity and regularization is central to the effectiveness of model finetuning. The BOFT framework enlarges the parameter space available during finetuning while introducing structured priors, facilitating an improved balance between expressivity and regularization.

\textbf{Spectral properties}. Orthogonal finetuning tends to produce superior spectral characteristics compared to LoRA, as it exactly maintains the spectral norm of the original weight matrix $\bm{W}^0$. This property becomes apparent by considering the singular value decomposition: $\bm{W}^0=\bm{U}\bm{\Sigma}\bm{V}^\top$, where $\bm{U}$ and $\bm{V}$ are orthogonal matrices, and $\bm{\Sigma}$ is diagonal with the singular values. In both OFT and BOFT, an orthogonal matrix $\bm{R}$ is applied from the left, resulting in the updated parameterization $\bm{R}\bm{U}\bm{\Sigma}\bm{V}^\top$, which leaves the maximal singular value (that is, the spectral norm of $\bm{W}^0$) unchanged. This invariance in spectral norm is known to enhance both training stability and the model's ability to generalize~\cite{miyato2018spectral,yoshida2017spectral}. Additional mathematical aspects are detailed in Appendix~\ref{app:sn_property}.

\textbf{Orthogonal finetuning as learning bilinear similarity}. The BOFT method can be formulated as learning a bilinear similarity of the form $\bm{w}^0_i\bm{R}\bm{x}$, with $\bm{w}^0_i$ being the $i$-th column (corresponding to a specific neuron) of $\bm{W}^0$. Here, BOFT can be interpreted as optimizing the bilinear similarity matrix $\bm{R}$, subject to the significant constraint of orthogonality, which naturally aligns with topics in distance metric learning~\cite{xing2002distance} and bilinear forms~\cite{roman2005advanced}.

\textbf{Inductive bias and generalization in BOFT}. Because $\bm{R}(m,b)$ in BOFT generally refers to a structured, lower-dimensional segment of the full orthogonal group, the approach effectively reduces the model’s hypothesis space, imparting a particular inductive bias. We posit that the explicit structure imposed by butterfly factorization fosters beneficial inductive biases for generalization. Such structure shares patterns found in traditional linear transforms, such as the discrete Fourier, sine/cosine, and Hadamard transforms~\cite{dao2019kaleidoscope}. Additionally, the use of sparse factorization within BOFT may further contribute to advantageous, implicit biases~\cite{gunasekar2017implicit,li2018algorithmic,liu2019neural}.

\textbf{Comparison to butterfly-based sparse training}. Numerous prior studies~\cite{dao2019learning,dao2019kaleidoscope,chen2022pixelated,dao2022monarch} have explored sparse training by leveraging the butterfly parameterization. These works predominantly center on reparameterizing neural network weight matrices with the butterfly structure and perform training from the initial randomization. In contrast, \cite{dao2022monarch} explores a finetuning approach wherein pretrained weights are first projected onto a specific form of butterfly matrices, after which the projected elements are adjusted for new downstream objectives. Distinct from these methods, BOFT introduces an alternative finetuning scheme that applies transformations to weights through shared matrices across layers.

\input{rebuttal-results}

\section{Concluding Remarks and Limitations}

This work introduces Orthogonal Butterfly, a broadly applicable parameter-efficient finetuning approach for foundation models that leverages the butterfly architecture. The principal idea for enhancing parameter-efficiency lies in expressing a dense orthogonal matrix as the product of several sparse orthogonal matrices. To facilitate the construction of such factorizations, a graphical information transmission framework is proposed. Utilizing this framework, we establish that the butterfly structure meets the desired properties for sparse orthogonal matrix decompositions. Our empirical results demonstrate the efficacy of BOFT in finetuning across a range of domains, including large language models, major vision architectures, and text-to-image generative systems. These findings also highlight the broad applicability and effectiveness of BOFT as a general finetuning mechanism.

While BOFT shows strong empirical performance, it is not without limitations. In particular, since the ultimate orthogonal matrix in BOFT results from composing multiple orthogonal matrices, training with BOFT incurs a marginally greater computational overhead compared to OFT. Reducing the training time for BOFT remains an unresolved question. Encouragingly, after finetuning concludes, the orthogonal matrices acquired via BOFT can be fused directly into the base model’s weights, eliminating any extra inference overhead. Furthermore, it is still uncertain whether the butterfly network constitutes the most optimal architecture for information propagation. Our proposed framework opens avenues for leveraging insights from the field of computer networking, where topological efficiency for information flow is a central focus. We anticipate that alternative, potentially more efficient network configurations could be exploited for the construction of dense orthogonal matrices.

\end{document}